\chapter{CƠ SỞ LÝ THUYẾT VÀ CÔNG TRÌNH LIÊN QUAN}
\label{ch:co_so}

\section{Bài toán phát hiện đối tượng}

\subsection{Phát biểu bài toán}

Cho một ảnh $I \in \mathbb{R}^{H \times W \times 3}$, bài toán phát hiện đối tượng yêu cầu sinh ra
một tập hợp
\begin{equation}
\label{eq:phatbieu}
\mathcal{D} = \{(c_i, b_i, s_i)\}_{i=1}^{N},
\end{equation}
trong đó $c_i \in \{1,\dots,K\}$ là nhãn lớp, $b_i = (x_i, y_i, w_i, h_i)$ là hộp bao và
$s_i \in [0,1]$ là điểm tin cậy. Số lượng vật thể $N$ không biết trước và thay đổi theo từng ảnh
--- đây chính là điểm khiến bài toán khó hơn phân lớp ảnh, nơi đầu ra luôn là một véc-tơ có kích
thước cố định.

Lịch sử phát triển của lĩnh vực đi qua hai dòng lớn. Dòng hai giai đoạn, tiêu biểu là Faster
R-CNN~\cite{ren2017fasterrcnn}, tách bài toán thành bước đề xuất vùng rồi mới phân loại và tinh
chỉnh hộp bao; độ chính xác cao nhưng chi phí tính toán lớn. Dòng một giai đoạn, khởi đầu từ
YOLO~\cite{redmon2016yolo} và SSD~\cite{liu2016ssd}, dự đoán trực tiếp hộp bao và lớp trên bản đồ
đặc trưng chỉ trong một lượt truyền xuôi, đổi một phần độ chính xác lấy tốc độ đủ cho ứng dụng thời
gian thực. Đề tài này làm việc hoàn toàn trong dòng thứ hai.

\subsection{Định dạng chú thích}

Hai định dạng chú thích được dùng trong đề tài. Định dạng YOLO lưu mỗi vật thể trên một dòng văn
bản gồm \texttt{class\_id}, tâm hộp $(x_c, y_c)$ và kích thước $(w, h)$, tất cả đã chuẩn hoá về
đoạn $[0,1]$ theo kích thước ảnh; đây là định dạng mà thư viện huấn luyện đọc trực tiếp. Định dạng
COCO~\cite{lin2014coco} dùng một tệp JSON duy nhất chứa ba khối \texttt{images},
\texttt{annotations} và \texttt{categories}, với hộp bao lưu theo góc trên trái và kích thước tuyệt
đối tính bằng điểm ảnh.

Phép chuyển từ YOLO sang COCO là bước bắt buộc trong đề tài vì khâu đánh giá dùng pycocotools. Công
thức chuyển đổi:
\begin{align}
\label{eq:yolo2coco}
x_{\text{trái}} &= \left(x_c - \tfrac{w}{2}\right) \cdot W_{\text{ảnh}}, &
y_{\text{trên}} &= \left(y_c - \tfrac{h}{2}\right) \cdot H_{\text{ảnh}}, \\
w_{\text{tuyệt đối}} &= w \cdot W_{\text{ảnh}}, &
h_{\text{tuyệt đối}} &= h \cdot H_{\text{ảnh}}. \nonumber
\end{align}
Trường \texttt{area} bằng tích $w_{\text{tuyệt đối}} \cdot h_{\text{tuyệt đối}}$. Trường này không
chỉ mang tính mô tả: bộ đánh giá COCO dùng chính nó để phân nhóm vật thể thành nhỏ, trung bình và
lớn khi tính $AP$ theo kích thước --- chỉ số trung tâm của đề tài này.

\section{Bộ phát hiện một giai đoạn}

\subsection{Nguyên lý chung}

Bộ phát hiện một giai đoạn bỏ hẳn bước đề xuất vùng. Ảnh đi qua một mạng xương sống tích chập để
sinh bản đồ đặc trưng, sau đó một đầu dự đoán gắn trực tiếp lên bản đồ đó sẽ xuất ra hộp bao và
điểm phân lớp cho từng vị trí. Vì toàn bộ quá trình chỉ cần một lượt truyền xuôi, tốc độ suy luận
cao hơn đáng kể so với dòng hai giai đoạn~\cite{redmon2016yolo}.

Điểm yếu cố hữu của cách làm này là mất cân bằng cực đoan giữa số vị trí nền và số vị trí chứa vật
thể. Focal Loss~\cite{lin2017focal} giải quyết bằng cách giảm trọng số đóng góp của các mẫu dễ:
\begin{equation}
\label{eq:focal}
\mathcal{L}_{\text{focal}}(p_t) = -\alpha_t \left(1 - p_t\right)^{\gamma} \log(p_t),
\end{equation}
trong đó $p_t$ là xác suất mô hình gán cho lớp đúng và $\gamma$ điều khiển mức độ giảm trọng số.

\subsection{Đa tỉ lệ và vật thể nhỏ}
\label{sec:vat_the_nho}

Một biển báo cách xe 50 mét và một biển báo cách 5 mét có kích thước trên ảnh chênh nhau hàng chục
lần. Mạng kim tự tháp đặc trưng (FPN)~\cite{lin2017fpn} xử lý điều này bằng cách kết hợp đặc trưng
ở nhiều tầng phân giải: tầng nông giữ chi tiết không gian tốt cho vật thể nhỏ, tầng sâu giàu ngữ
nghĩa hơn cho vật thể lớn.

Dù vậy, vật thể nhỏ vẫn là điểm yếu dai dẳng~\cite{tong2020smallobject}. Nguyên nhân nằm ở bước
sải: đầu dự đoán P3 hoạt động ở bước sải 8, nghĩa là một vật thể 16 điểm ảnh chỉ trải trên hai ô
lưới. Tín hiệu mà một vật thể nhỏ để lại trên bản đồ đặc trưng vốn đã yếu, nên nó là thành phần
\emph{mong manh nhất} của toàn bộ hệ thống --- và cũng là thành phần đầu tiên bị tổn hại khi ta can
thiệp vào mô hình bằng các phép nén. Đây chính là lý do lý thuyết khiến đề tài đo $AP$ tách theo
nhóm kích thước thay vì chỉ đọc mAP tổng.

\subsection{Hàm mất mát của họ YOLO hiện đại}

Các phiên bản YOLO hiện đại~\cite{jocher2023ultralytics} là bộ phát hiện không dùng hộp neo. Hàm
mất mát tổng gồm ba thành phần:
\begin{equation}
\label{eq:yololoss}
\mathcal{L} = \lambda_{\text{box}} \mathcal{L}_{\text{CIoU}}
            + \lambda_{\text{cls}} \mathcal{L}_{\text{BCE}}
            + \lambda_{\text{dfl}} \mathcal{L}_{\text{DFL}}.
\end{equation}

Thành phần thứ nhất dùng CIoU~\cite{zheng2020diou}, một mở rộng của IoU có tính thêm khoảng cách
giữa hai tâm hộp và độ lệch tỉ lệ khung, nhờ đó vẫn cung cấp gradient hữu ích ngay cả khi hai hộp
chưa giao nhau --- trường hợp mà IoU thuần tuý cho gradient bằng không. Thành phần thứ ba là
Distribution Focal Loss~\cite{li2020gfl}, coi mỗi cạnh hộp bao là một phân phối rời rạc thay vì một
giá trị điểm, giúp việc hồi quy biên ổn định hơn khi ranh giới vật thể mờ.

Khi bật chưng cất tri thức, một thành phần thứ tư được cộng thêm vào biểu
thức~\eqref{eq:yololoss} --- chi tiết ở Mục~\ref{sec:kd_lythuyet}.

Sau khi suy luận, bước hậu xử lý NMS loại bỏ các hộp trùng cho cùng một vật thể: các hộp được sắp
theo điểm tin cậy giảm dần, hộp có điểm cao nhất được giữ lại và mọi hộp còn lại có IoU với nó vượt
ngưỡng sẽ bị loại.

\section{Hệ thống độ đo}

\subsection{IoU, Precision và Recall}

Độ đo nền tảng là tỉ số giao trên hợp giữa hộp dự đoán $B_p$ và hộp thật $B_g$:
\begin{equation}
\label{eq:iou}
\text{IoU}(B_p, B_g) = \frac{|B_p \cap B_g|}{|B_p \cup B_g|} .
\end{equation}
Một dự đoán được tính là dương tính đúng khi lớp khớp và IoU vượt ngưỡng cho trước. Từ đó suy ra
\begin{equation}
\label{eq:pr}
\text{Precision} = \frac{TP}{TP + FP}, \qquad
\text{Recall} = \frac{TP}{TP + FN}.
\end{equation}

\subsection{mAP}

Sắp toàn bộ dự đoán theo điểm tin cậy giảm dần và tính cặp $(\text{Precision}, \text{Recall})$ tại
từng ngưỡng cho ta đường cong PR. Độ chính xác trung bình $AP$ là diện tích dưới đường cong này, và
$mAP$ là trung bình của $AP$ trên toàn bộ các lớp. Cách định nghĩa này kế thừa từ bộ chuẩn PASCAL
VOC~\cite{everingham2010pascalvoc} và được COCO~\cite{lin2014coco} mở rộng thêm chiều ngưỡng IoU:
\begin{equation}
\label{eq:map}
AP = \int_{0}^{1} p(r)\, dr, \qquad
mAP = \frac{1}{K} \sum_{k=1}^{K} AP_k .
\end{equation}
Quy ước COCO báo cáo hai biến thể: mAP@0,5 tính tại một ngưỡng IoU duy nhất là 0,5, và
mAP@0,5:0,95 lấy trung bình trên mười ngưỡng từ 0,5 đến 0,95 với bước 0,05. Biến thể thứ hai khắt
khe hơn nhiều vì nó đòi hỏi hộp bao phải khít, chứ không chỉ chồng lấn ở mức tối thiểu.

Một hệ quả quan trọng của việc mAP là độ đo \emph{xếp hạng}: khi tính mAP, không được lọc dự đoán
theo ngưỡng tin cậy trước, vì làm vậy sẽ cắt mất phần đuôi của đường cong PR và bóp méo kết quả.
Ngưỡng tin cậy chỉ nên áp dụng ở khâu hiển thị.

\subsection{$AP$ tách theo nhóm kích thước}
\label{sec:ap_theo_kichthuoc}

Đây là công cụ trung tâm của đề tài, nên cần nói rõ vì sao nó cần thiết.

Biểu thức~\eqref{eq:map} lấy trung bình trên các lớp nhưng \emph{không} phân biệt kích thước vật
thể. Nếu tập đánh giá có nhiều vật thể lớn --- vốn dễ phát hiện --- thì $AP$ cao trên nhóm này sẽ
kéo mAP tổng lên và che lấp một suy giảm nghiêm trọng ở nhóm nhỏ. Hiện tượng pha loãng này càng
nặng khi phân bố kích thước lệch, đúng như trường hợp bộ dữ liệu được dùng ở đây (xem
Mục~\ref{sec:thanhphan}).

Bộ đánh giá COCO~\cite{lin2014coco} khắc phục bằng cách tính $AP$ riêng cho ba nhóm dựa trên diện
tích tuyệt đối của hộp bao:
\begin{equation}
\label{eq:apsize}
AP_{\text{nhỏ}}: a < 32^2, \qquad
AP_{\text{trung bình}}: 32^2 \le a < 96^2, \qquad
AP_{\text{lớn}}: a \ge 96^2,
\end{equation}
với $a$ tính bằng điểm ảnh vuông. Ba chỉ số này cho phép trả lời câu hỏi mà mAP tổng không trả lời
được: khi một phép can thiệp làm mô hình kém đi, \emph{phần kém đó rơi vào nhóm nào}. Với bài toán
biển báo giao thông, câu trả lời có ý nghĩa an toàn trực tiếp, vì biển ở xa --- tức vật thể nhỏ ---
mới là thứ quyết định thời gian phản ứng của xe.

\subsection{Độ đo chi phí triển khai}

Ba đại lượng bổ sung được dùng để phản ánh ràng buộc thực tế: độ trễ suy luận tính bằng mili-giây
cho mỗi ảnh, tốc độ khung hình bằng nghịch đảo của độ trễ, và dung lượng mô hình tính bằng
megabyte. Ngoài giá trị trung vị, đề tài còn báo cáo phân vị 95\% của độ trễ, vì với hệ thống thời
gian thực thì trường hợp xấu nhất mới là ràng buộc thật chứ không phải trường hợp trung bình.

Cả ba đại lượng đều phụ thuộc phần cứng nên phải luôn báo cáo kèm thiết bị đo. Ngoài ra, do GPU
thực thi bất đồng bộ, phép đo thời gian phải có bước đồng bộ hoá tường minh và một giai đoạn khởi
động trước khi bấm giờ, nếu không sẽ ghi nhận sai thời gian tính toán thực.

\section{Nén mô hình}

\subsection{Chưng cất tri thức}
\label{sec:kd_lythuyet}

Chưng cất tri thức~\cite{hinton2015distillation} huấn luyện một mô hình học trò nhỏ bắt chước phân
phối đầu ra mềm của một mô hình thầy lớn hơn. Hàm mất mát kết hợp
\begin{equation}
\label{eq:distill}
\mathcal{L} = (1-\alpha)\,\mathcal{L}_{\text{cứng}}
            + \alpha\, T^{2}\, \mathcal{L}_{\text{KL}}\!\left(\sigma(z_s/T) \,\|\, \sigma(z_t/T)\right),
\end{equation}
với $T$ là nhiệt độ làm mềm phân phối và $z_s, z_t$ là logit của học trò và thầy.

Trong bối cảnh phát hiện đối tượng, việc áp dụng chưng cất phức tạp hơn phân lớp ảnh vì đầu ra gồm
cả thành phần hồi quy hộp bao lẫn thành phần phân loại, và vì mất cân bằng nền--vật thể rất lớn.
Các công trình chuyên biệt đã đề xuất chưng cất trên đặc trưng trung gian thay vì chỉ trên
logit~\cite{chen2017kddetection,li2017mimicking}; một khảo sát tổng quan về các biến thể có thể tìm
thấy trong~\cite{gou2021kdsurvey}.

Giả thuyết nền là phân phối mềm của thầy mang cái thường được gọi là ``tri thức tối'': thông tin về
\emph{quan hệ} giữa các lớp mà nhãn cứng không có. Chẳng hạn khi nhìn một biển ``Speed Limit 30'',
thầy có thể gán xác suất thứ hai cho ``Speed Limit 80'' --- hàm ý rằng hai lớp này dễ nhầm --- trong
khi nhãn cứng chỉ nói lớp đúng là gì.

Tuy nhiên lợi ích này \emph{không được bảo đảm}, và điều kiện áp dụng thường bị bỏ qua trong các
trình bày phổ thông. Khi khoảng cách năng lực giữa thầy và trò quá nhỏ, lượng tri thức tối có thể
truyền đạt là rất ít, trong khi chi phí nhiễu do bắt chước một phân phối không hoàn hảo vẫn phải
trả đủ. Khi bản thân thầy chưa đủ tốt, học trò đang bắt chước một tín hiệu sai. Và khi trọng số
$\alpha$ quá lớn, thành phần bắt chước có thể lấn át tín hiệu từ nhãn cứng. Đề tài này quan sát
được đúng một tình huống như vậy, trình bày ở Chương~\ref{ch:ket_qua}.

Trong bối cảnh phát hiện đối tượng, thành phần chưng cất được cộng thêm vào hàm mất mát tổng ở
biểu thức~\eqref{eq:yololoss} và xuất hiện trong nhật ký huấn luyện dưới dạng một cột riêng --- chi
tiết này về sau trở thành bằng chứng để xác minh cơ chế đã thực sự được kích hoạt.

\subsection{Lượng tử hoá}
\label{sec:quant_lythuyet}

Lượng tử hoá thay biểu diễn dấu phẩy động 32 bit bằng biểu diễn có độ rộng bit nhỏ hơn. Phép ánh xạ
affine sang số nguyên 8 bit~\cite{jacob2018quantization} có dạng
\begin{equation}
\label{eq:quant}
r \approx S\,(q - Z),
\end{equation}
trong đó $r$ là giá trị thực, $q$ là giá trị nguyên, $S$ là hệ số tỉ lệ và $Z$ là điểm không. Các
tham số $S$ và $Z$ được xác định ở bước \emph{hiệu chuẩn}: chạy mô hình trên một tập ảnh đại diện để
quan sát dải giá trị thực tế của các kích hoạt. Việc chọn dải này ảnh hưởng đáng kể tới chất lượng
mô hình sau lượng tử hoá, và các hướng dẫn thực hành chi tiết có thể tìm thấy
trong~\cite{krishnamoorthi2018quantizing,nagel2021quantwhitepaper}; một khảo sát tổng quan về các
phương pháp lượng tử hoá được trình bày trong~\cite{gholami2022quantsurvey}.

Lượng tử hoá sau huấn luyện hấp dẫn vì không cần huấn luyện lại. Cái giá là sai số lượng tử hoá:
mỗi giá trị thực bị làm tròn về mức gần nhất trong 256 mức khả dĩ. Ở mức FP16, việc giảm độ rộng
bit gần như không ảnh hưởng tới chất lượng mô hình trong nhiều tác
vụ~\cite{micikevicius2018mixedprecision}. Ở mức INT8, mức nén mạnh hơn nhiều nhưng sai số cũng lớn
hơn hẳn.

Có một dự đoán lý thuyết đáng chú ý ở đây, và nó được kiểm chứng ở Chương~\ref{ch:ket_qua}: vì sai
số lượng tử hoá tác động lên toàn bộ bản đồ đặc trưng, còn vật thể nhỏ vốn chỉ để lại tín hiệu yếu
trên bản đồ đó (Mục~\ref{sec:vat_the_nho}), \emph{tỉ lệ tín hiệu trên nhiễu của vật thể nhỏ suy
giảm nhiều hơn của vật thể lớn}. Nói cách khác, chi phí của lượng tử hoá được dự đoán là phân bố
\emph{không đều} theo kích thước vật thể --- và đó chính là thứ mAP tổng không thể hiện được.

Ngoài ra, lợi ích về tốc độ của INT8 không tự động: nó chỉ hiện thực hoá khi môi trường thực thi có
nhân tính toán số nguyên được tối ưu cho phần cứng đích. Nếu không, thời gian tiết kiệm được có thể
bị tiêu hết vào các phép chuyển đổi lượng tử hoá và giải lượng tử hoá xen giữa các lớp.

\section{Công trình liên quan về biển báo giao thông}

GTSRB~\cite{stallkamp2012gtsrb} là bộ chuẩn kinh điển cho \emph{nhận dạng} biển báo, gồm các ảnh đã
cắt sát quanh từng biển; đây là bài toán phân lớp chứ không phải phát hiện.
GTSDB~\cite{houben2013gtsdb} bổ sung phần \emph{phát hiện} với ảnh đường phố đầy đủ.
TT100K~\cite{zhu2016tt100k} mở rộng quy mô sang bối cảnh Trung Quốc với ảnh toàn cảnh độ phân giải
cao, còn Mapillary Traffic Sign~\cite{ertler2020mapillary} cung cấp dữ liệu đa quốc gia thu từ
camera hành trình.

Phân biệt giữa bộ dữ liệu kiểu nhận dạng và bộ dữ liệu kiểu phát hiện có ý nghĩa trực tiếp với đề
tài này: như sẽ trình bày ở Mục~\ref{sec:thanhphan}, bộ dữ liệu được dùng là một tập trộn chứa cả
hai kiểu, và chính điều đó chi phối phân bố kích thước vật thể --- yếu tố quyết định cách đọc mọi
chỉ số về sau. Hiện tượng bộ dữ liệu mang thiên lệch đặc thù và mô hình học phải thiên lệch đó thay
vì học đặc trưng tổng quát đã được chỉ ra từ lâu~\cite{torralba2011bias}, còn việc chỉ số trên tập
kiểm tra không chuyển được sang dữ liệu mới cũng đã được ghi nhận trên các bộ chuẩn
lớn~\cite{recht2019imagenetv2}. Riêng vấn đề trùng lặp gần giữa các tập, vốn làm chỉ số bị thổi
phồng một cách âm thầm, đã được khảo sát kỹ trên nhiều bộ chuẩn phổ biến~\cite{barz2020duplicates}.
